\documentclass{report}
\usepackage{graphicx}
\usepackage{listings}
%\usepackage{fullpage}

\begin{document}

% \documentclass[12pt]{article}

% \usepackage{ngerman}
% \selectlanguage{ngerman}
% \usepackage[latin1]{inputenc}

% \usepackage{graphicx}

% \begin{document}

\pagestyle{empty}

\begin{center}

\vspace*{-2.8cm}
\begin{minipage}[c]{.30\textwidth}
  \begin{flushleft}
    \includegraphics[width=3cm,clip=]{LOGO/new-aerlogo.eps}%
  \end{flushleft}
\end{minipage}
\begin{minipage}[c]{.43\textwidth}
    { ~\\ Lehrstuhl f\"{u}r Aerodynamik \\ Prof.~Dr.-Ing. N.~A.~Adams}%
\end{minipage}
\begin{minipage}[c]{.25\textwidth}
  \begin{flushright}
    \vspace*{1em}
    \includegraphics[width=3cm,clip=]{LOGO/TUMLogo_oZ_Vollfl_blau_RGB.eps}%
  \end{flushright}
\end{minipage}

\vspace*{3.3cm}
\begin{minipage}[c]{11cm}
{\LARGE\bf 
A Riemann solver }
\end{minipage}

\vspace*{0.8cm}
Mark F\"{o}rster\\

\vspace*{2.8cm}
{\bfseries Diplomarbeit}

\vspace*{1.2cm}
%\large
\normalsize
\vfill
\begin{tabular}{ll}
Betreuer: &  Dr.-Ing. Christian Stemmer \\
\ & \ \\
Ausgabe: & 1.\ Oktober 2008\\
Abgabe: & 31.\ M\"{a}rz 2009\\
\end{tabular}

\vspace*{1.6cm}
%\large
Lehrstuhl f\"{u}r Aerodynamik an der Technischen Universit\"{a}t M\"{u}nchen\\
2009
\end{center}

\pagebreak
\pagestyle{plain}

%\end{document}



\tableofcontents
\chapter{Introduction}
\label{sec:intro}


This work aimes at the numerical simluation of high enthalpy flow in
porosities. The subject is of interest, as CITATION
NEEDED (Prof Hornung) has discovered porous surfaces to inhibit/block a
certain form of instabilities in hypersonic flows over a body/cone. A numerical
simulation of this configuration is desired to furnish data for a detailed
analysis of the acting phenomena. The simulation method of choice for the first approach
to this configuration is called Smoothed Particle Hydrodynamics (SPH).Although
there has been no known application of SPH to this precise/concrete case yet/before,
SPH is adopted, as it promises a crucial advantage in the way boundary
conditions are implemented. And for the treatment of porosities,
boundaries/walls represent a non negligible part of the problem.

SPH designates a meshfree method for numerical simulation of particle
ensembles (respectively matter that can be modelled as an ensemble of
particles as it is the case for fluids).
Thereby,the term particle ensembles not only refers
to fluids (as the name of the method would suggest) but also to solids. In
fact nowadays SPH is employed to simulate problems in different fields,
such as body elasticity and fracture,liquid flows and gas
dynamics~\cite{Monaghan2005}.
The initial developement of the method dates back to 1977, where Gingold and
Monaghan tempted to compute an astrophysical problem with a new, more
efficient method~\cite{Gingold1977}. Independently and simultaneously, work in
this field was also carried out by Lucy who came up with the same
idea~\cite{Lucy1977}.
ADVANTAGES???DISADVANTAGES???

PREVIOUS APPLICATIONS IN GAS DYNAMICS (citing different relevant papers)

OTHER APPROACHES FOR CAVITY/POROSITY SIMULATION (look up in literature)





\chapter{Method}
\label{sec:method}
\section{Basic Idea of SPH}
\label{sec:BasicsSPH}

To illustrate the basic idea of SPH, it is helpful to first have a look at the
ruling equations of fluid dynamics, here represented by the SPH equations in
lagrangian form. The use of the lagrangian form seems logical as SPH is a
particle method.

\begin {equation}
NAVIER STOKES IN LAGRANGIAN FORM
\end {equation}

The above equations have the form~\cite{Monaghan2005}

\begin {equation}
\label{eq:EFD_form}
{dA(r)\over dt}=f(A(r),\nabla A(r),r)
\end {equation}
where
\begin {equation}
{d(.)\over dt}={\partial(.)\over \partial t}+v\nabla(.)
\end{equation}

Thus the rates of change of  physical quantity depend on its spatial
derivatives. The goal of any method for numerical simulation is to approximate
these derivatives by information of a discrete number of points in order for a
computer to be able to handle it. 


The SPH method is based on the idea of approximating a function $A(x)$ by an
integral interpolant
\begin{equation}
\label{eq:int_interpolant}
A_I(x)=\int A(r')W(r-r',h)dr'
\end{equation}
where $W(r-r',h)$ is the so called kernel or smoothing function and $h$ the
smoothing length. This Interpolation is exact, if $W$ is the delta function,
otherwise it describes/represents an approximation. As the theoretically
perfect Kernel, the delta function, is of no practical use, the Kernels
employed in practice are functions which tend to the delta function, when $h$
tends to zero. Examples for those kernel functions are a Gaussian Kernel or
Kernels based on Schoenbergs~\cite{Schoenberg1946} $M_n$ splines~\cite{Monaghan2005}. More details on Kernels are given in the section
(LINK TO SECTION??) and in the book on SPH published by Liu~\cite{Liu2003}

The fact of using a kernel different from the delta function constitutes/is the
first of two approximations /approximative assumptions made for the SPH
method. It is often refered to as kernel approximation~\cite{Liu2003}.
The second approximation consists of a discretization of the integral
expression in equation\ref{eq:int_interpolant}, leading to a summation over
all discretized elements (the particles) of the domain instead
of an integration when calculating the interpolant value $A_I(x)$, which
finally allows a treatment of the problem by a numerical machine. In SPH
literature, this second approximation is commonly called particle
approximation~\cite{Liu2003}

Applying the particle approximation on \ref{eq:int_interpolant} leads with the
relation between mass, density and volume for each discretized particle
\begin{equation}
m_i=\rho_i V_i
\end{equation}
 to the following expression

\begin{equation}
A_S(r)=\sum_b m_b \frac{A_b}{\rho_b}W(r-r_b,h)
\end{equation}

In the same way, the derivative of the fieldfunction $frac{dA(x)}{dx}$ can be
approximated to give\cite{Monaghan2005}or\cite{Liu2003}(in Liu with gradient,
im Monaghan for 1 direction)

\begin{equation}
\label{eq:simpleDerivative}
\nabla A_S(r)=\sum_b m_b \frac{A_b}{\rho_b}\nabla W(r-r_b,h)
\end{equation}

The disadvantage of the above equation \ref{eq:simpleDerivative} is, that for
$A(r)=const.$ the approximated derivative does not vanish. By writing (in
Monaghan the equation figures not with gradient but only in 1 direction!?!)
\begin{equation}
\label{eq:derivativeIdentity}
\nabla A = \frac{1}{\Phi}(\nabla {\Phi A}-A\nabla \Phi)
\end{equation}
where $\Phi$ is any differentiable function, this problem can be resolved. In its
SPH form  equation \ref{eq:derivativeIdentity} reads as follows:
\begin{equation}
\nabla_a A = \frac{1}{\Phi_a}\sum_b frac{m_b}{\rho_b}(A_b-A_a)\nabla_a W_{ab}
\end{equation}


It can be seen, that the (approximation of the) derivative of the
fieldfunction $A$ is now expressed by a summation over the fieldfunction itself
and the derivative of the kernel function - which can be determined exactly.

Now the basic formulations are introduced in principle: Both the fieldfunction
itself and its spacial derivatives can be approximated using only information
(=the states) of a discrete number of particles. Hence, the change rates of
the states described by the equations of fluid dynamics, which have a form
according to \ref{eq:EFD_form} can be rewritten using this information. Doing
this one obtains the SPH formulation of the equations of fluid dynamics.

The detailed derivation of those equations will not be shown here, as the
interest  only lies in describing the principle ideas behind the SPH
method. To resume, this are the following ones: A (field) function can be
described as an integral interpolation over the function itself and a kernel
function (which has to be chosen in an appropriate way). A further
approximation consists in transforming the integral over the domain into a
summation over the discrete elements constituing this domain (the particles). This allows to express both
the field function and its spacial derivatives dependent exclusively on discrete values
of the fieldfunction itself (and the kernel function, which is known as well as
its derivatvies).By finally choosing an appropriate kernel function with
compact support, the summation over all particles in the domain can be reduced
to a summation over those particles within the support domain of the particle
in question.

As already mentionned above, the SPH formulation of the equations of fluid
dynamics is obtained by applying the above principles and conducting different
transformations to the original (exact) equations. Depending on the
transformations used, the final SPH equations can assume a different form. In
the following, a common form for the change rates of density, velocity and
internal energy based on the EULER EQUATIONS (i.e. without viscosity) is
listed\cite{Monaghan2005}\cite{Liu2003} (perhaps list NSG-form as well, and
explain why Euler is of interest as well (for 1DSPH code, perhaps later real
code including viscosity???)

\begin{equation}
\frac{d\rho _a}{\mathit{dt}}=\sum{m_{b}v_{\mathit{ab}}\nabla _{a}W_{\mathit{ab}}}
\end{equation}

\begin{equation}
\label{eq:VelChangeRate}
\frac{dv_{a}}{\mathit{dt}}=-\sum {m_{b}(\frac{P_{a}}{\rho_{a}^{2}}+\frac{P_{b}}{\rho _{b}^{2}}+\Pi _{ab})\nabla_{a}W_{ab}}
\end{equation}


\begin{equation}
\frac{de_{a}}{\mathit{dt}}=-\mathit{}\frac{1}{2}\sum{m_{b}(\frac{P_{a}}{\rho _{a}^{2}}+\frac{P_{b}}{\rho _{b}^{2}}+\Pi_{\mathit{ab}})v_{\mathit{ab}}\nabla _{a}W_{\mathit{ab}}}
\end{equation}



For the detailed procedure of derivation and for alternative formulations of
these equations refer to \cite{Monaghan2005}\cite{Liu2003}

Artificial Viscosity

Another poit that has to be taken into account for SPH simulations, especially
such including shock phenomena, is arificial viscosity\cite{Monaghan2005}. It
prevents particles from penetrating one into an other in zones with high
compression rates and also generally stabilizes a numerical
allgorithm.
Artificial viscosity for numerical simulation of shock phenomena
in general was first introduced by \cite{vonNeumann1950}. They suggested to
introduce an artificial dissipation to give the shock a thickness which
allowed it to be calculated numerically (in the cited example with finite difference
scheme) without introduction of pre-/after shock boundary conditions.    ...(see article)
The first use of artificial viscosity specifically for the mesh free SPH goes
back to Lucy\cite{Lucy1977}.He suggested an artificial bulk viscosity to
prevent a slow buildup of acoustic energy by integration errors in zones with
low density respectively low particle presence. 
Finally, the form of artificial viscosity most commomly used in SPH
literature\cite{Liu2003} is the one proposed by Monaghan and
Gingold\cite{Monaghan1983}. They stated that the use of one of the artificial
viscosities cited above (Neumann Richtmyer (adapted to SPH) or bulk
viscosity) on shock problems would lead to either excessive particle
oscillation or excessive smearing of the shock front. The expression for this
artificial viscosity $\Pi_{ab}$, which is added to the pressure term of
equation \ref{eq:VelChangeRate}, is:

\begin{equation}
\Pi _{\mathit{ab}}=-\nu
(\frac{v_{\mathit{ab}}r_{\mathit{ab}}}{r_{\mathit{ab}}^{2}+\epsilon
h_{\mathit{ab}}^{2}})
\end{equation}
where 
\begin{equation}
\nu =\frac{h_{\mathit{ab}}}{\rho _{\mathit{ab}}}(\alpha
c_{\mathit{ab}}-\beta
\frac{h_{\mathit{ab}}v_{\mathit{ab}}r_{\mathit{ab}}}{r_{\mathit{ab}}^{2}+\epsilon
h_{\mathit{ab}}^{2}})
\end{equation}

$v_{ab}$ stands for $v_a-v_b$, $r_{ab}$ for $r_a-r_b$, $c_{ab}$ is the
averaged sound speed of the two particles in question $1/2(c_a+c_b)$. The same
applies for $\rho_{ab}(=1/2(\rho_a+\rho_b))$ and $h_{ab}=(1/2(h_a+h_b))$. IT
HAS TO BE MENTIONED SOMEWHERE (but probably not right here) THAT: the use of
the averaged smoothing length $h_{ab}$ is a way to symmetrize particle
interactions when h is spatially variable. In this case a particle $a$ could be
located in the support domain of a particle $b$ but not vice versa. This would
lead to an assymetric mutual contribution of forces. To avoid this, using the
averaged smoothing length between two particles is an appropriate remedy\cite{Liu2003}.


 
mention somewhere that kernel unity is inverse of volume?!?!?

\section{Implementation of an SPH code}
This section presents some particular features of an SPH code. Emphasis is
placed on these features which are finally implemented in one of the
two SPH codes designed and modified during this work.
Points of interest are the implementation of different boundary conditions in
SPH, a closer description of the kernel function, as well as different methods/allgorithms
to perform the search for the interacting particles. Furthermore different
equations of state are introduced, depending on the nature of the flow problem
to be considered(compressible/incompressible). Finally some words on numerical
integration schemes that have proven to be reliable for SPH applications.


\subsection{boundary condition implementation}
For particles at the edge of the calculation domain, the summation interpolation does not
give exact results, as the integral (or better the sum) is truncated due to
the simple fact that outside the domain there are no
particles.\cite{Liu2003}(find better citation...) This efect is called the
problem of article deficiencythe particle deficiency
problem.\cite{Liu2003}. In some cases there is no need to consider this
problem, as the areas close to the edges are not of interest and due to
limited simulation time there is no time for the error to propagate into the
actual zone of interest. The 1D shock tube simulation described further below
(PERHAPS LINK TO SECTION)is an example for such a case.  
But for other applications the edges play an important role or the domain has
to be rendered periodocally. So there is a need to define specific boundary
conditions. This can be accomplished by a special boundary treatment, which
often includes so called virtual particles or ghost
particles\cite{Liu2003}. The first ones to come up with this idea were
\cite{LIBERSKY1993}. Virtual particles are used to model solid walls by
placing them inside the solid body, mirrored at the surface. For each real
particle whose support domain interfres with the wall, one such particle is
built. Virtual particles have the same properties as their real counterpart
except for the velocity. Depending on the type of boundary condition they are
supposed to represent, different velocities are assigned to them.\cite{Hu2006}

non slip velocity boundary condition

\begin{equation}
v_{virtual}=2v_{wall}-v_{real}
\end{equation}

free slip velocity boundary condition

\begin{equation}
v_{virtual}^{tangential}=v_{real}^{tangential}
\end{equation}

ILLUSTRATION OF VELOCITY ASSIGNMENT FOR VIRTUAL BOUNDARY PARTICLES

For the treatment of periodic boundary conditions, no ghost particles are
needed. Particles leaving the calculation domain on one side are simple
reinserted on the opposite side (with the same properties).


\subsection{kernel function}
???not sure, if its the right place here, but somewhere there should be a
paragraph with more detailed information on kernels used, compact support,
perhaps even some properties,...(and perhaps some graphics???)


Some properties of Kernel functions (smoothing functions) and notions related
to them have allready been mentionned in the section
\ref{sec:BasicsSPH}. In principal every function can act as Kernel function
when it respects a set of criteria complsed by \cite{Liu2003}.
The most important ones are:

The smoothing function should vanish outside of a certain area around its
center. This property is called compact support and the domain where $W\neq0$
the support domain of the smoothing function. The smoothing length $\kappa h$ and the
support length $h$ may not be mixed up.
\begin{equation}
W(x-x')=0,\text{for}|x-x'|>\kappa h
\end{equation}

The smoothing function should be normalized to 1, which ensures that constants
are interpolated exactly
\begin{equation}
\int{W(x-x',h)dx'}=1
\end{equation}

The smoothing function should tend towards the delta function as the smoothing
length $h$ tends to zero
\begin{equation}
\lim\limits_{h \rightarrow 0}{W(x-x',h)}=\delta(x-x')
\end{equation}
This is to ensure consistency of the integral approximation, which means that
the approximative expression tends to the exact expression of the integral
interpolant equation.

The full list of SPH kernel properties can be found in  \cite{Liu2003}.
Based on those properties a variety of Kernels are used in SPH literature. A
very commonly used Kernel though is the $M_4$ Kernel based on Schoenbergs
$M_n$-Splines\cite{Schoenberg1946}.

FORMULA M_4 Kernel

In the work of \cite{Fulk1996} a mesure of merit for Kernels is developed and
they came to the result, that in 1D no kernel performed significantly better
than the cubic spline.
This is why for the 1D SPH code described in section (WHATEVER) the cubic
splime kernel is employed.

\subsection{neighboring particle search methods}

\subsection{equation of state}

\subsection{numerical integration}
...conservation properties more important than order of scheme\cite{Monaghan2005}



\section{1D SPH code}
\subsection{short overview}
\subsection{components}
\subsection {structure}
\section{2D SPH code}
\subsection{short overview}
\subsection{components}

\subsection{structure}

\lstinputlisting[language=c++,caption=An example of program listing,captionpos=b]{code/hello.cpp}

\chapter{Results}
\section{1D SPH code}
...Monaghan05 p.1739... not width important but pre and after shock values...

\section{...}
\label{sec:results}

In the Figure~\ref{fig:1dshock} one can see \ldots

\begin{figure}[h]
  \centering
  \includegraphics[width=0.95\textwidth]{img/1dshock.eps}
  \caption{An example of a figure}
  \label{fig:1dshock}
\end{figure}





\chapter{Conclusion}
\label{sec:conclusion}

\listoffigures

\bibliography{bibdata}
\bibliographystyle{plain}

\end{document}

